% ============================================================
%  深圳大学实验报告（SZU Experiment Report Template）
%  编译方式：XeLaTeX（fontspec + xeCJK）
% ============================================================

\documentclass[a4paper,12pt]{article}

% ----------------------------
% 字体：中文 Noto，西文 Times 系
% ----------------------------
\usepackage{fontspec}
\usepackage{xeCJK}
\IfFontExistsTF{Noto Serif CJK SC}
  {\setCJKmainfont{Noto Serif CJK SC}}
  {\setCJKmainfont{SimSun}}
\IfFontExistsTF{Times New Roman}
  {\setmainfont{Times New Roman}}
  {\setmainfont{TeX Gyre Termes}}

\usepackage{geometry}
\geometry{left=1in, right=1in, top=1in, bottom=1in}

\usepackage{longtable}
\usepackage{graphicx}
\usepackage{fancyhdr}
\usepackage{tikz}
\usetikzlibrary{calc}
\usepackage{verbatim}
\usepackage{float}
\usepackage{array}
\usepackage{url}

\fancyhf{}
\fancyhead[C]{}
\fancyhead[R]{}

\setlength{\parskip}{2pt}
\setlength{\emergencystretch}{3em}
\renewcommand{\figurename}{图}
\renewcommand{\tablename}{表}

\begin{document}

% ============================================================
% 封面
% ============================================================
\begin{titlepage}
    \centering
    \vspace*{1.6cm}
    \Huge{\textbf{深圳大学实验报告}}\\[1.4cm]

    {\Large\textbf{课程名称}\quad\underline{\makebox[8cm][c]{Python程序设计}}}\\[0.5cm]
{\Large
\noindent\makebox[3.2cm][l]{\textbf{项目名称}}
\underline{\makebox[8cm][c]{小项目 1 NumPy 图像编码器}}}\\[0.5cm]
{\Large
\noindent\makebox[3.2cm][l]{\textbf{学 \quad \quad 院}}
\underline{\makebox[8cm][c]{人工智能学院}}}\\[0.5cm]
{\Large
\noindent\makebox[3.2cm][l]{\textbf{专 \quad \quad 业}}
\underline{\makebox[8cm][c]{人工智能}}}\\[0.5cm]
{\Large
\noindent\makebox[3.2cm][l]{\textbf{指导教师}}
\underline{\makebox[8cm][c]{舒婷}}}\\[0.5cm]
{\Large
\noindent\makebox[3.2cm][l]{\textbf{报 \quad 告 \quad 人}}
\underline{\makebox[3cm][c]{李云舟}}
\hspace{0.5cm}
\textbf{学号}
\underline{\makebox[3cm][c]{2025681110}}}\\[0.5cm]
{\Large
\noindent\makebox[3.2cm][l]{\textbf{实验时间}}
\underline{\makebox[8cm][c]{2026 年 11 月 24 日至 25 日}}}\\[0.5cm]
{\Large
\noindent\makebox[3.2cm][l]{\textbf{提交时间}}
\underline{\makebox[8cm][c]{2026 年 11 月 26 日}}}\\[1.2cm]

    \vfill
    \Large{\textbf{教务处制}}
\end{titlepage}

\newpage

% ============================================================
% 正文
% ============================================================

\section{Objectives}

本实验用 NumPy 手工复现 OpenCLIP ViT-B/32 的图像塔，不加载 PyTorch 或 OpenCLIP 包，也不更新任何参数。输入是课程预处理后的 $[B,3,224,224]$ 浮点数组（最短边缩放到 224、中心裁剪、按 OpenCLIP 的均值与方差逐通道归一化），输出是每张图一个 512 维单位向量，需要与课程提供的独立 FP32 参考在 $\mathrm{atol}=2\times10^{-5}$、$\mathrm{rtol}=2\times10^{-4}$ 内一致，中间 token 的容差为 $2\times10^{-4}$。

实现只允许使用 reshape、transpose、矩阵乘与广播；课程提供的 LayerNorm、精确 GELU、注意力系数与单位化函数保持原样，\texttt{check.py}、\texttt{run.py} 与预处理脚本也不得改动。除数值一致外，还要说明三个问题：张量形状如何逐级变化、多头注意力的四个轴如何组织、以及调试时如何定位最早出现偏差的阶段。代码托管于 \url{https://github.com/Lyzz-11/szu-python-project1-vit}，README 记录了环境、依赖与运行步骤。

\section{Overview}

实验在课程附件给出的 10 张测试图、152 个视觉张量（约 326 MB）与逐阶段参考数组上完成。模型规模固定为：宽度 $D=768$、头数 $H=12$、每头 $D/H=64$、12 层、MLP 中间宽度 3072、token 数 50（1 个 CLS 加 49 个图块）、输出 512 维，全程 float32。参数约定有三条必须遵守：线性层权重按 $[\mathrm{out},\mathrm{in}]$ 存放，使用时需转置；\texttt{visual.proj} 已经是 $[\mathrm{in},\mathrm{out}]$，不能转置；patch 投影没有偏置，图像块也没有 LayerScale。

整体流程是：图像切成 $7\times7$ 个 $32\times32$ 图块并展平，经无偏置线性投影后与 CLS、可学习位置编码拼成 50 个 token；随后通过 12 个 pre-LN 块（多头注意力、两项残差、GELU MLP），取 CLS 行做 post-Norm，投影到 512 维并单位化。表~\ref{tab:flow} 记录了各步的形状。

\begin{table}[H]
\centering
\caption{从输入图像到输出向量的形状变化（$B$ 为批大小）}
\label{tab:flow}
\small
\begin{tabular}{|l|l|l|}
\hline
步骤 & 主要操作 & 形状 \\
\hline
输入 & 预处理后的 RGB 张量 & $[B,3,224,224]$ \\
切块 & reshape、transpose 后展平 & $[B,49,3072]$ \\
嵌入 & 线性投影、拼 CLS、加位置编码 & $[B,50,768]$ \\
12 层编码 & 注意力与 MLP，宽度 768 / 3072 & $[B,50,768]$ \\
输出头 & 取 CLS、post-Norm、$768\rightarrow512$、单位化 & $[B,512]$ \\
\hline
\end{tabular}
\end{table}

\section{Implementations}

\textbf{切块与嵌入。}\texttt{patchify} 先把空间轴拆成两段，得到 $[B,3,7,32,7,32]$，再用 \texttt{transpose} 把轴排成（批、图块行、图块列、通道、块内行、块内列），最后合并成 $[B,49,3072]$：

\begin{verbatim}
split   = images.reshape(B, 3, 7, 32, 7, 32)
ordered = split.transpose(0, 2, 4, 1, 3, 5)
patches = ordered.reshape(B, 49, 3072)
\end{verbatim}

投影使用 $W_{\mathrm{patch}}^{\top}$，不加偏置；CLS 按批复制后拼在 token 轴最前，再加上同一个位置编码。

\textbf{多头注意力。}线性层权重按 $[\mathrm{out},\mathrm{in}]$ 存放，因此 $qkv=xW^{\top}+b$；特征轴 reshape 成 $(H,D/H)=(12,64)$ 后，把头轴提到 token 轴之前，使每个头独立做一次 $N\times N$ 的注意力：
\[
S=\frac{QK^{\top}}{\sqrt{64}},\qquad A=\mathrm{softmax}(S),\qquad O=(AV)W_o^{\top}+b_o .
\]
其中 $K^{\top}$ 只交换最后两个轴；合并头时按相反顺序 \texttt{transpose} 再 \texttt{reshape} 回 $[B,N,768]$，最后乘输出投影。注意力系数直接调用课程提供的稳定实现（先减行最大值再取指数），图像塔不启用 causal 掩码，但该开关保留给最终项目的文本塔复用。

\textbf{残差块与 MLP。}每个块先做 $\mathrm{LN}_1$ 与注意力并相加得到 $r$；$\mathrm{LN}_2$ 之后把宽度拓到 3072，过精确 erf 版 GELU，压回 768 再做第二次相加：
\[
r=x+\mathrm{Attn}\!\left(\mathrm{LN}_1(x)\right),\qquad
y=r+\mathrm{MLP}\!\left(\mathrm{LN}_2(r)\right).
\]

\textbf{输出头。}取 \texttt{tokens[:,0,:]} 做 post-Norm，乘以已是 $[\mathrm{in},\mathrm{out}]$ 的 \texttt{visual.proj}，最后调用课程的单位化函数，得到 $[B,512]$。

\section{Results}

\textbf{复现环境与运行。}Windows 11、Python 3.12.14、NumPy 2.3.5、Pillow 12.3.0，纯 CPU，不需要 GPU 与联网。在 \url{project1-numpy-student} 目录下执行：

\begin{verbatim}
python -m pip install -r requirements.txt
python check.py --stage all
python run.py --output output_tokens.npz
\end{verbatim}

\textbf{分阶段数值校验。}\texttt{check.py} 依次比对六个阶段的输出，全部通过，实测误差见表~\ref{tab:stage}。判定条件是形状、dtype、有限性与容差同时满足，容差按阶段区分：输出端 $2\times10^{-5}$，内部 token $2\times10^{-4}$。

\begin{table}[H]
\centering
\caption{各阶段输出与参考实现的最大、平均绝对误差（rtol 均为 $2\times10^{-4}$）}
\label{tab:stage}
\small
\setlength{\tabcolsep}{3.5pt}
\begin{tabular}{|l|l|c|c|c|c|}
\hline
阶段 & 输出形状 & 最大绝对误差 & 平均绝对误差 & 容差 & 结果 \\
\hline
preprocess   & $(10,3,224,224)$ & $0$ & $0$ & atol $2\times10^{-6}$ & 通过 \\
patches      & $(2,49,3072)$    & $0$ & $0$ & atol $2\times10^{-4}$ & 通过 \\
tokens       & $(2,50,768)$     & $4.387\times10^{-5}$ & $1.335\times10^{-6}$ & atol $2\times10^{-4}$ & 通过 \\
block0       & $(2,50,768)$     & $7.147\times10^{-6}$ & $3.348\times10^{-7}$ & atol $2\times10^{-4}$ & 通过 \\
last\_tokens & $(2,50,768)$     & $4.196\times10^{-5}$ & $2.633\times10^{-6}$ & atol $2\times10^{-4}$ & 通过 \\
output       & $(10,512)$       & $5.662\times10^{-7}$ & $8.688\times10^{-8}$ & atol $2\times10^{-5}$ & 通过 \\
batch        & $(3,512)$        & $4.061\times10^{-7}$ & $7.848\times10^{-8}$ & atol $2\times10^{-5}$ & 通过 \\
\hline
\end{tabular}
\end{table}

\textbf{关键输出。}\texttt{run.py} 在 10 张图上写出 $[10,512]$ 的 float32 结果，图像 ID 顺序与 \url{assets/images.json} 一致；每行的 $\ell_2$ 范数落在 $[0.99999988,1.0]$，满足单位长度。用 $[1,0,1]$ 的乱序重复批重算，与逐张推理的最大偏差为 $4.1\times10^{-7}$，说明实现既不依赖批内位置，也没有把整批当作全局统计量。10 张图两两余弦相似度最低 0.4624、平均 0.5721、最高 0.6739；\texttt{image\_01} 的前五个分量为 $-0.0399,\ 0.0850,\ 0.0799,\ 0.0214,\ 0.0210$，最大单元素绝对值约 $0.196$。整轮 \texttt{--stage all} 约 12 秒，其中课程提供的逐元素 GELU 占绝大部分时间。

\textbf{运行开销与误差趋势。}权重加载与编码器初始化约 $1.7$ 秒（$326$ MB 的 npz 解压为主）；单张前向 $0.39$ 秒，十张 $3.9$ 秒，基本随批大小线性增长，说明瓶颈在逐元素 GELU 而非矩阵乘。误差随层数的变化并不单调：\texttt{block0} 为 $7.1\times10^{-6}$，到 \texttt{last\_tokens} 累积到 $4.2\times10^{-5}$，而 \texttt{output} 反而回落到 $5.7\times10^{-7}$，原因是输出端的 post-Norm 与单位化把幅值归一到 1 附近，相对误差被同步压缩。

\section{Discussion}

\textbf{形状变化与一处顺序错误。}切块最初写成一次 \texttt{reshape(B,49,3072)}，形状正确但元素乱序：NumPy 按行优先展平，每一段 3072 会跨越通道与整行像素。必须先把空间轴拆成两段 $[B,3,7,32,7,32]$，再 \texttt{transpose} 成（批、图块行、图块列、通道、块内行、块内列），最后才合并。前两步只改变轴的解读方式，第三步才真正重排内存，所以顺序由 \texttt{transpose} 决定而非 \texttt{reshape}——这也是 \texttt{patches} 阶段误差恰好为 $0$ 的原因。

\textbf{一处真正的 bug：4 维张量上的转置。}注意力最初写成 \texttt{scores = q @ k.T}。对二维矩阵 \texttt{.T} 就是转置，但对 $[B,H,N,64]$ 的 $K$，\texttt{.T} 会把四个轴全部反转，得到 $[64,N,H,B]$，把 batch 与 head 一并打乱，\texttt{check.py --stage block0} 直接报核心维度不匹配（$12$ 与 $64$）。改成只交换最后两个轴的 \texttt{k.transpose(0,1,3,2)} 后，\texttt{block0} 误差降到 $7.1\times10^{-6}$。这次排查还暴露了一个更隐蔽的问题：先跑的 \texttt{--stage tokens} 只验证 \texttt{embed}，根本不经过注意力，带 bug 的版本在该阶段依然通过，必须跑到 \texttt{block0} 才能确认注意力正确。因此“阶段通过”只说明该阶段覆盖的那部分正确，定位错误要先确认检查的覆盖范围。

\textbf{结果分析。}预处理与切块的误差为 $0$，说明张量搬运没有引入任何数值改动；中间层误差在 $10^{-5}$ 量级，来自 BLAS 的累加顺序与参考实现不同，属于浮点实现的正常差异；输出端因为做了 post-Norm 与单位化，幅值归一后相对误差被压到 $10^{-7}$，所以最严的容差反而余量最大。所有阶段都保留了一个数量级以上的余量，说明实现没有隐藏的偏差累积。

\textbf{局限与可改进之处。}本实现只覆盖视觉塔，不含文本塔、tokenizer 与图文相似度；权重固定，不含训练或微调逻辑，改 \texttt{config.json} 不会变成另一个模型；实现假定输入已完成预处理，直接送原图会与参考不一致。性能瓶颈是课程提供的逐元素 GELU（$B=10$ 时 12 个块约 $1.8\times10^{7}$ 次 erf 调用），若允许可用向量化实现替代，但 tanh 近似会改变数值、无法满足当前容差；12 个块之间严格串行，进一步提速只能靠批量矩阵乘或降低精度。接入最终项目时，应把 \texttt{image\_encoder.py} 与 \texttt{helpers.py} 一并复制到对应目录。

\textbf{本人贡献与使用的帮助。}本人完成六个 TODO 的实现、逐阶段验收与全部数据整理，未修改课程提供的 \texttt{helpers.py}、\texttt{preprocess\_image.py}、\texttt{runtime.py} 与 \texttt{check.py}；过程中使用 DeepSeek Harness 辅助排查 4 维转置问题并整理表格与文字。文中数字均取自本机 \texttt{check.py} 的实际输出，代码已逐行复核。

\newpage

% ============================================================
% 批阅与成绩评定页
% ============================================================
\begin{tikzpicture}[remember picture, overlay]
  \draw[thick]
    ($(current page.north west)+(2cm,-3cm)$)
    rectangle
    ($(current page.south east)+(-2cm,3.5cm)$);
\end{tikzpicture}

\vspace{1cm}

\noindent \textbf{指导教师批阅意见：}
\vspace{6cm}
\hfill

\vspace{1cm}

\noindent \textbf{成绩评定：}
\vspace{2cm}
\hfill

\vspace{1cm}

\noindent
\hfill
\textbf{指导教师签字：}\quad\quad\quad\quad\quad\\[0.5cm]

\hfill
\textbf{年}\quad\quad
\textbf{月}\quad\quad
\textbf{日}\quad

\vspace{1cm}

\noindent \textbf{备注：}

\begin{tikzpicture}[remember picture, overlay]
  \node[
    anchor=north west,
    inner sep=0pt,
    text width=\dimexpr\paperwidth-4cm\relax
  ] at ($(current page.south west)+(2cm,3.1cm)$) {
    \begin{minipage}{\dimexpr\paperwidth-4cm\relax}
      \small
      \noindent
      \textbf{注：}\;1、报告内的项目或内容设置，可根据实际情况加以调整和补充。\\[0.2cm]
      \phantom{\textbf{注：}\;}2、教师批改学生实验报告时间应在学生提交实验报告时间后 10 日内。
    \end{minipage}
  };
\end{tikzpicture}

\end{document}
